\chapter{Stage VIII：Regge--Wheeler 方程式と black-hole QNM}

\solproblem{VIII.1　Schwarzschild tortoise 座標}{
Schwarzschild metric から $r_*$ と両端の漸近形を導け。}{
$f(r)=1-2M/r$ とし radial null line $0=-f\dd t^2+f^{-1}\dd r^2$ を
$t\mp r_*=\mathrm{const.}$ にするには $\dd r_*/\dd r=f^{-1}$。従って
\begin{equation}
 r_*=r+2M\log\left(\frac r{2M}-1\right)+C.
\end{equation}
$C=0$ とすると horizon $r=2M+\epsilon$ で
\begin{equation}
 r_*=2M\log\frac{\epsilon}{2M}+2M+O(\epsilon)\to-\infty,
\end{equation}
infinity で
$r_*=r+2M\log(r/2M)+O(M^2/r)\to+\infty$。従って外部領域は
$r_*\in\real$ 全体へ写る。}

\solproblem{VIII.2　spherical harmonics の parity と odd metric 成分}{
scalar/vector/tensor harmonic の parity を示し odd sector の自由成分を数えよ。}{
parity map $(\vartheta,\varphi)\mapsto(\pi-\vartheta,\varphi+\pi)$ で
$Y_{lm}\mapsto(-1)^lY_{lm}$。polar vector $D_AY$ は $(-1)^l$、sphere volume
formを掛けた axial vector $X_A=-\varepsilon_A{}^BD_BY$ は追加の負号を得て
$(-1)^{l+1}$。polar/axial tensor もそれぞれ $(-1)^l,(-1)^{l+1}$。
fixed $(l,m)$ の odd metric は
$h_tX_A$、$h_rX_A$、$h_2X_{AB}$ の3 scalar amplitude。odd gauge function
$\xi X_A$ 一つで $h_2=0$ にでき、2 amplitude が残る。Einstein constraint
が一つを nondynamical にし、最終的な propagating master field は1。
$l=0$ に odd harmonic はなく、$l=1$ で $X_{AB}=0$ なので別扱い。}

\solproblem{VIII.3　linearized Einstein 方程式から covariant Regge--Wheeler 方程式へ}{
warped-product identities を用いた projection と elimination を示せ。}{
$\rmd s^2=g_{ab}\dd x^a\dd x^b+r^2\Omega_{AB}\dd\theta^A\dd\theta^B$ とし、
odd perturbation を $h_{aB}=h_aX_B$、$h_{AB}=h_2X_{AB}$ と展開する。
gauge invariant combination は
\begin{equation}
 \widetilde h_a=h_a-\frac12\nabla_ah_2+\frac{r_a}{r}h_2.
\end{equation}
$D^AX_A=0$、$D^2X_A=[1-l(l+1)]X_A$、
$D^BX_{AB}=-(l-1)(l+2)X_A/2$ を linearized Ricci tensor へ代入すると
\begin{align}
 0=P_a={}&-\Box_2\widetilde h_a+\nabla_a\nabla_b\widetilde h^b
 +\frac2r(r_b\nabla_a\widetilde h^b-r_a\nabla_b\widetilde h^b)\\
 &-\frac2{r^2}r_ar_b\widetilde h^b+\frac{l(l+1)}{r^2}\widetilde h_a,\\
 0=P={}&\nabla_a\widetilde h^a.
\end{align}
$\lambda=(l-1)(l+2)$ とし
\begin{equation}
 \vsub{\Psi}{odd}=\frac{2r}{\lambda}\varepsilon^{ab}
 \left(\nabla_a\widetilde h_b-\frac2r r_a\widetilde h_b\right)
\end{equation}
を定義する。$\varepsilon^{ab}\nabla_aP_b$ を取り、背景 identity
$\nabla_a\nabla_br=(M/r^2)g_{ab}$、$r_ar^a=f$、constraint $P=0$ を使うと
\begin{equation}
 \varepsilon^{ab}\nabla_aP_b=-\frac{\lambda}{2r}
 \left[\Box_2-\frac{l(l+1)}{r^2}+\frac{6M}{r^3}\right]\vsub{\Psi}{odd}.
\end{equation}
従って bracket が0。Bianchi identity は残りの projected equation が従属
であることを保証する。これが成分を飛ばさない gauge-invariant 導出である。}

\solproblem{VIII.4　Regge--Wheeler potential と低 multipole}{
potential を導き両端を検査し、$l=0,1$ の例外を説明せよ。}{
$\Box_2=-f^{-1}\partial_t^2+\partial_r(f\partial_r)$ と
$\partial_{r_*}=f\partial_r$ を前問へ代入し
\begin{equation}
 [-\partial_t^2+\partial_{r_*}^2-\vsup{V_l}{RW}(r)]\vsub{\Psi}{odd}=0,
 \quad \vsup{V_l}{RW}=f\left[\frac{l(l+1)}{r^2}-\frac{6M}{r^3}\right]
\end{equation}
を得る。horizon では $f\to0$ なので $V\to0$ exponentially in $r_*$、
infinity では $V=l(l+1)/r^2+O(r^{-3})\to0$。$l=0$ には odd vector harmonic
が存在しない。$l=1$ では $\lambda=0$ で tensor harmonic と radiative
master definition が消え、vacuum static solution は background の angular
momentumを変える infinitesimal Kerr perturbation である。従って wave
potential formula を機械的に $l=0,1$ へ適用しない。}

\solproblem{VIII.5　horizon と infinity の符号}{
$\rme^{-\rmi\omega t}$ 規約で QNM 条件を vanishing connection coefficient として書け。}{
horizon の future-regular ingoing wave は advanced time $v=t+r_*$ の関数
$\rme^{-\rmi\omega v}$ なので radial factor は
$\rme^{-\rmi\omega r_*}$。infinity の outgoing wave は retarded time
$u=t-r_*$ の関数で radial factor $\rme^{+\rmi\omega r_*}$。horizon-normalized
解を $\vsup{f_H}{in}$ とし infinity で
\begin{equation}
 \vsup{f_H}{in}=\vsub{\mathcal C}{out}(\omega)
 \rme^{+\rmi\omega r_*}+\Cincoming(\omega)
 \rme^{-\rmi\omega r_*}
\end{equation}
と展開する。QNM 条件は $\Cincoming(\omega)=0$。$\im\omega<0$ では両端の
QNM radial factor は実 $r_*$ 上で発散するが、これは符号誤りでなく
resonance boundary condition である。}

\solproblem{VIII.6　barrier-top WKB による fundamental QNM}{
Schwarzschild fundamental mode を leading WKB で見積もり $M$ scaling と比較せよ。}{
frequency equation $[-\partial_{r_*}^2+V]\psi=\omega^2\psi$ の barrier
maximum $r_0$ で $V\simeq V_0+V_0''(r_*-r_{*0})^2/2$ とすると
\begin{equation}
 \omega^2\simeq V_0-
\rmi(n+1/2)\sqrt{-2V_0''}.
\end{equation}
gravitational $l=2$ では
$r_0/M=(9+\sqrt{17})/4=3.280776\ldots$、
$M^2V_0=0.1512867$、$M^4V_0''=-0.00991941$。$n=0$ より
\begin{equation}
 M\vsup{\omega_{20}}{WKB}\simeq0.39885-0.08829\rmi.
\end{equation}
高精度値 $M\omega_{20}\simeq0.37367-0.08896\rmi$ と比べ、damping は良く
実部は leading order の誤差を示す。metric を $r=M\bar r$、
$r_*=M\bar r_*$ と scale すれば方程式は $(M\omega)^2$ のみを含むので
$\omega\propto M^{-1}$。}

\solproblem{VIII.7　tortoise coordinate の complex scaling}{
QNM exposure 条件と inverse map の branch による制限を導け。}{
両端で $r_*\mapsto\rme^{\rmi\theta}r_*$ とする。右端 outgoing wave と左端
ingoing wave（$r_*=-s$）はいずれも magnitude
\begin{equation}
 \left|\rme^{\rmi\omega s\rme^{\rmi\theta}}\right|
 =\rme^{-s\im(\omega\rme^{\rmi\theta})}
\end{equation}
を持つので $\im(\omega\rme^{\rmi\theta})>0$、すなわち
$\theta>-\arg\omega$ が exposure 条件。ただし
\begin{equation}
 r=2M\left[1+W\left(\rme^{r_*/(2M)-1}\right)\right]
\end{equation}
で inverse map は Lambert $W$ の branch を含む。global straight rotation
は logarithm/W branch cut、$r=0$ singularity、他 sheet の horizon image を
横切り得る。従って実際には potential が解析な有限内部区間を保ち、両 tail
だけを曲げる exterior/smooth complex scaling を用いる。}

\solproblem{VIII.8　Gaussian basis 行列と conditioning}{
Regge--Wheeler 方程式の overlap/kinetic 行列を作り精度喪失を診断せよ。}{
一例として中心 $x_j$ と range $b_j$ を持つ
$\phi_j(r_*)=(r_*-x_j)^{p_j}\rme^{-\alpha_j(r_*-x_j)^2}$ を使う。
$p=0$、共通中心なら
\begin{equation}
 N_{ij}=\sqrt{\frac\pi{\alpha_i+\alpha_j}},\qquad
 T_{ij}=\frac{2\alpha_i\alpha_j\sqrt\pi}
 {(\alpha_i+\alpha_j)^{3/2}},\qquad
 V_{ij}=\int\phi_i\vsub{V}{RW}\phi_j\dd r_*.
\end{equation}
complex scaling 後は derivative に $\rme^{-2\rmi\theta}$、potential argument
に deformation を入れる。geometric range ratio $q$ が1に近すぎると columns
がほぼ従属、遠すぎると中間 scale が欠落する。polynomial order 増大も
moment 行列の dynamic range を増やす。
\begin{checkbox}
$N$ の smallest eigenvalue と condition number、Cholesky/SVD cutoff、
generalized residual、左右 overlap、working precision を記録する。precision
を倍にして安定桁が増えない eigenvalue は物理 pole と認定しない。
\end{checkbox}}

\solproblem{VIII.9　Schwarzschild--de Sitter の二 horizon と Nariai limit}{
二 horizon を $r_*\to\pm\infty$ へ写し near-Nariai barrier scaling を求めよ。}{
$f=1-2M/r-\Lambda r^2/3$ の正根を $r_b<r_c$ とする。各単純根で
$f\simeq2\kappa_b(r-r_b)$、
$f\simeq2\kappa_c(r_c-r)$、
$\kappa_h=|f'(r_h)|/2$。従って
\begin{align}
 r_*&\sim\frac1{2\kappa_b}\log(r-r_b)\to-\infty,\\
 r_*&\sim-\frac1{2\kappa_c}\log(r_c-r)\to+\infty.
\end{align}
$\delta=\sqrt{1-9\Lambda M^2}\downarrow0$ と置くと
\begin{align}
 r_{b,c}&=3M\left(1\mp\frac{\delta}{\sqrt3}+O(\delta^2)\right),\\
 r_c-r_b&=2\sqrt3M\delta+O(\delta^2),\\
 \kappa_{b,c}&\sim\frac{\delta}{3\sqrt3M}.
\end{align}
coordinate $r$ では static region が縮むが $r_*$ 幅は
$\order(\kappa^{-1})\sim M/\delta$ に広がる。master potential の高さは
$\order(\kappa^2)$、従って near-Nariai QNM frequency も $\order(\kappa)$。}

\solproblem{VIII.10　QNM pole と continuum level density}{
QNM pole contribution と reference-subtracted cut contribution を比較せよ。}{
適切な source $F$ と reference operator に対し
\begin{equation}
 \Delta_F(E)=-\frac1\pi\im\left[
 \bra{F}\Resolvent(E+\rmi0)\ket{F}-
 \bra{F}\vsub{\Resolvent}{ref}(E+\rmi0)\ket{F}\right].
\end{equation}
contour を変形すると単純 QNM $z_R=\omega_R^2$ の項は
\begin{equation}
 \Delta_R(E)=-\frac1\pi\im\frac{Z_F}{E-z_R},\qquad
 Z_F=\frac{\mathcal N_F(z_R)}{\partial_E\Cincoming(z_R)},
\end{equation}
残りが reference-subtracted branch cut と他 poles。$Z_F$ は複素で cut は
energy dependent なので、全 response の extremum 条件
$\partial_E[\Delta_R+\vsub{\Delta}{cut}]=0$ は一般に
$E=\re z_R$ で満たされない。さらに $E=\omega^2$ の非線形写像により
frequency plot の peak も $\re\omega_R$ と一致しない。pole position、residue、
full real-axis peak を別々に報告する。}
